\documentclass[11pt]{article}

\usepackage[margin=1in]{geometry}
\usepackage{amsmath,amssymb,amsthm}
\usepackage{mathtools}
\usepackage{enumitem}
\usepackage{hyperref}
\usepackage{natbib}

\newtheorem{proposition}{Proposition}
\newtheorem{corollary}{Corollary}
\newtheorem{definition}{Definition}
\newtheorem{assumption}{Assumption}
\newtheorem{remark}{Remark}
\newtheorem{lemma}{Lemma}

\title{Strategic Under-Specification: A Stackelberg Game Between User and Assistant}
\author{Pranjal Agarwal \\ \small BITS Pilani, Hyderabad Campus}
\date{}

\begin{document}
\maketitle

\begin{abstract}
Users of coding assistants routinely under-specify their intent, and the assistant's
policy for handling missing information---guess silently, or ask a clarifying
question---is itself a strategic choice that users respond to. We formalize this as
a Stackelberg game: the assistant commits to an ask/guess policy, and a population
of users with privately known, heterogeneous specification costs best-respond. We
first solve a decision-theoretic benchmark in which the policy is exogenous. Here,
improving the assistant's raw guessing accuracy alone cannot make users worse
off, even though it induces them to specify less---an immediate envelope-theorem
consequence of letting a single agent optimize freely against an improved
technology, but one worth stating explicitly because it is the benchmark against
which the paper's actual moral-hazard channel (a gap between believed and true
policy) should be measured; naive ``better guessing causes worse specs and worse
outcomes'' stories conflate the two. We then
close the loop and solve for the assistant's own optimal policy under two regimes:
a \emph{pooling} policy restricted to a single population-wide ask rate, and a
\emph{screening} policy that can condition on the user's declared specification. We
show that an unbiased (welfare-aligned) assistant achieves full efficiency under
screening despite users' private information---incentive compatibility is free when
principal and agent share an objective---so all screening-regime inefficiency is
attributable to misalignment between the assistant's effective objective and true
user welfare, not to private information per se. Pooling, by contrast, loses
efficiency from restricted instruments even when unbiased. We characterize both
equilibria in closed or near-closed form under a tractable linear-risk
specification, derive testable comparative statics, evaluate them via numerical
and LLM-based simulation experiments, and outline a proposed human-study design for
future work.
\end{abstract}

\section{Introduction}

A user typing a coding request faces a choice: spend effort articulating edge
cases, naming conventions, and constraints, or leave them implicit and trust the
assistant to infer them. The assistant, in turn, must decide for each unspecified
detail whether to guess or to ask. Neither choice is made in a vacuum: the user's
specification effort depends on what they expect the assistant to do with silence,
and the assistant's ask/guess policy is chosen (explicitly, or as an emergent
property of training) with some implicit model of how users will respond. Section 2
treats the assistant's policy as exogenous and solves the user's problem against
it---a one-sided optimization, useful as a benchmark but not yet a game. Sections
4--7 close the loop: the assistant is a Stackelberg leader that commits to a policy
anticipating the population's best response, and we solve for that policy under two
different instrument restrictions, then compare them.

\subsection{Related literature}

This work bridges contract theory and the emerging literature on formal models of human--AI interaction.

\paragraph{Contract theory and screening without transfers.}
Our screening framework builds on classical adverse-selection models~\citep{myerson1981,baron-myerson1982,mussa-rosen1978,laffont-tirole1986} that exploit the Spence--Mirrlees single-crossing property~\citep{spence1973} to separate heterogeneous agents. However, unlike classical monopoly or regulatory screening where the principal uses monetary transfers to extract information rents, here the assistant screens strictly through an interactive service margin (the ask rate $a$ offered in response to specification effort $m$). This connects to non-monetary screening and costly verification~\citep{townsend1979}, where information revelation is governed by operational friction rather than cash transfers.

\paragraph{Costly communication and strategic disclosure.}
The user's specification decision formalizes costly disclosure under strategic interpretation, connecting to strategic information transmission~\citep{crawford-sobel1982}, voluntary quality disclosure~\citep{grossman1981}, and rational inattention~\citep{sims2003}. In our setting, under-specification arises not from an incentive to deceive, but from the user's rational economizing on specification costs when anticipating automated assistance.

\paragraph{Stackelberg formulations of human--AI interaction.}
While game-theoretic formulations of hierarchical bi-level interactions date back to \citet{stackelberg1934}, human prompt specification in human--AI interaction has largely been examined through empirical studies in software engineering, where under-specified or implicit prompts frequently induce costly clarification and repair loops. To our knowledge, our model provides the first formal contract-theoretic microfoundation for this interaction, characterizing how assistant asking policies and guessing capabilities strategically shape user specification effort.

\section{Environment and the Decision-Theoretic Benchmark}
\label{sec:baseline}

\subsection{Setup}

A task has $k$ binary, independent, intent-relevant attributes
$\theta = (\theta_1,\dots,\theta_k) \in \{0,1\}^k$, each drawn i.i.d.\ uniform and
privately known to the user. The user chooses a specification effort level
$m \in \{0,\dots,k\}$ and truthfully reveals $m$ attributes.\footnote{We assume
truthful revelation conditional on the choice to specify: the strategic margin is
\emph{how much} to say, not whether to lie about what is said. A cheap-talk layer
on top of this costly-specification layer is a natural extension; see
Section~\ref{sec:discussion}.} Specifying $m$ attributes costs $c(m)$, with
$c(0)=0$, $c' > 0$, $c'' \ge 0$.

For each of the $k-m$ unspecified attributes, the assistant either \emph{asks}
(revealing the true value, at a cost to the user) or \emph{guesses}, succeeding
independently with probability $g \in (0,1)$. We model the assistant's policy at
this stage as a uniform ask rate $a \in [0,1]$ applied to every unspecified
attribute, so each is correctly resolved with probability
$q(a,g) = a + (1-a)g$. Task success requires \emph{all} $k$ attributes to be
correctly resolved (conjunctive: one wrong parameter or convention typically breaks
the deliverable). The user's payoff is
\[
U(m; a, g) \;=\; V \, q(a,g)^{\,k-m} \;-\; c(m) \;-\; a(k-m)\kappa ,
\]
with $V>0$ the value of success and $\kappa>0$ the user's cost of answering one
clarifying question. We treat $m$ as continuous on $[0,k]$ for tractability.

\subsection{Proposition 1: accuracy alone does not create moral hazard}

\begin{proposition}\label{prop:accuracy}
Fix $a$. Let $U^*(g) = \max_{m \in [0,k]} U(m;a,g)$. Then $U^*$ is weakly
increasing in $g$.
\end{proposition}

\begin{proof}
By the envelope theorem, at the optimum $m^*(g)$,
$dU^*/dg = \partial U/\partial g\big|_{m=m^*(g)} = V(k-m^*)q^{k-m^*-1}(1-a) \ge 0$.
The envelope theorem lets us ignore $dm^*/dg$: because $m^*$ is chosen optimally,
its first-order response to $g$ has no first-order effect on $U^*$.
\end{proof}

\begin{assumption}[Regularity]\label{ass:reg}
In the relevant range of $m$, $(k-m)\ln q(a,g) \ge -1$.
\end{assumption}

\begin{remark}[Interpreting Assumption~\ref{ass:reg}, and Monotonicity Survival]
\label{rmk:assumption1-audit}
The condition has no closed-form economic reading on its own; it is a sufficient
condition for the single-crossing property used in Corollary~\ref{cor:mono} and
Proposition~\ref{prop:bias} via Topkis's monotone comparative statics theorem~\citep{topkis1978,topkis1998,milgrom-shannon1994}, not a claim about primitives
directly. The direction is pinned down by the cross-partial itself: for fixed $a$,
\[
\frac{\partial^2 U}{\partial m\,\partial g} \;=\; -V(1-a)\,q^{\,k-m-1}\Big[(k-m)\ln q + 1\Big],
\]
so $U$ has \emph{decreasing differences} in $(m,g)$ --- the condition Topkis's
theorem~\citep{topkis1978,topkis1998} needs to conclude $m^*(g)$ is weakly decreasing --- exactly when
$(k-m)\ln q + 1 \ge 0$, i.e., when Assumption~\ref{ass:reg} holds as stated here.
If the inequality were reversed, $U$ would exhibit increasing differences and hence $m^*(g)$ would be locally increasing
(e.g., at $V{=}10,k{=}5,a{=}0,c(m){=}m^2/2$, $m^*$ rises from
$\approx0.26$ at $g{=}0.5$ to $\approx0.67$ at $g{=}0.9$ before turning over and
falling as $g\to1$). Loosely, since $\ln q \approx -(1-q)$ for $q$ near 1, the
condition is close to $(k-m)(1-q) \le 1$: the expected number of \emph{wrong}
attributes under the current policy is at most about one. A numerical grid evaluation over
$k\in[5,15]$, $q\in[0.70,0.95]$, and $\kappa\in[0.20,2.00]$ (726 grid points,
quadratic cost $c(m)=\kappa m^2/2$, evaluated at each point's own optimizing
$m^*$) finds that under the benchmark simulation valuation ($V=100$), the
condition holds on $89.9\%$ (653/726) of this parameter space.
In lower-valuation regimes where specification effort is lower
($V=10$, as in the numerical example above), the condition holds on roughly
$34\%$ (and $31.7\%$ at $V=8$); the exact share is sensitive to the scaling of
$c(m)$ relative to $V$. Crucially, across the tested parameter grid (2,552
parameter intervals), a comprehensive finite-difference evaluation found no violations
of monotonicity outside the sufficient condition: $m^*(g)$ remained weakly
decreasing in $g$ across all $2{,}162$ tested intervals where
Assumption~\ref{ass:reg} failed, providing numerical evidence that monotonicity
of specification can survive in practice even when the sufficient condition is
violated.
\end{remark}

\begin{corollary}\label{cor:mono}
Under Assumption~\ref{ass:reg}, $m^*(g)$ is weakly decreasing in $g$ (decreasing
differences in $(m,g)$; Topkis's theorem).
\end{corollary}

\begin{proof}
Immediate from the cross-partial computed in Remark~\ref{rmk:assumption1-audit}:
Assumption~\ref{ass:reg} is exactly the condition $\partial^2U/\partial m\,\partial g \le 0$
on the relevant range, which is decreasing differences in $(m,g)$; Topkis's
theorem then gives $m^*(g)$ weakly decreasing.
\end{proof}

So, holding the assistant's policy fixed, users specify less as the assistant gets
better at guessing, yet are never worse off for it. We flag explicitly that
Proposition~\ref{prop:accuracy} is a low bar: it is a direct application of the
envelope theorem to a single agent optimizing freely against a technology
parameter that enters its payoff with non-negative derivative, and would hold in
essentially any model of this form. Its purpose here is not to be a surprising
finding but to be the correct benchmark against which any claimed moral hazard
should be measured, so that Proposition~\ref{prop:bias} below can be read as
identifying the actual, non-trivial source of harm. It isolates the
question this paper is really about: what happens once the policy $a$ is no longer
exogenous, but chosen by a strategic assistant that anticipates this response?

\section{Sources of Strategic Under-Specification}
\label{sec:sources}

Proposition~\ref{prop:accuracy} assumed the assistant's ask rate is fixed and that
users correctly anticipate $q(a,g)$. Both assumptions can fail.

\subsection{Cost externality}
If a third party (a reviewer, downstream production systems) bears part of the
cost of a wrong guess, the user's private $(\kappa, V)$ understate the social
values, and the user under-specifies relative to the social optimum at any fixed
$g$---worsening, by Corollary~\ref{cor:mono}, as $g$ rises. Standard; not our focus.

\subsection{Miscalibrated belief updating under silent failure}
If wrong guesses are not always detected, and detection probability $d(m)$ is
increasing in $m$, users infer $g$ from experience rather than observing it, and
under-specifying users are also less likely to catch their own errors---an
upward-biased $\hat g$ that reinforces under-specification. This is a dynamic
learning story rather than a static equilibrium result; we return to it in the
empirical design (prediction P5).

\subsection{Belief about the assistant's policy}
\label{sec:naive-soph}
Let users be \emph{sophisticated} (correctly infer the assistant's true ask rate
$a_A$) or \emph{naive} (believe the assistant follows the welfare-optimal rate
$a^\dagger \ge a_A$).

\begin{proposition}\label{prop:bias}
Under Assumption~\ref{ass:reg}, naive users choose
$m^*_{\text{naive}} \le m^*_{\text{soph}}$ and
$U(m^*_{\text{naive}};a_A,g) < U(m^*_{\text{soph}};a_A,g)$: naive users
under-specify and are strictly worse off than sophisticated users, holding raw
guessing accuracy $g$ fixed throughout.
\end{proposition}

\begin{proof}
Because $q(a,g)=a+(1-a)g$ is strictly increasing in $a$ for $g\in(0,1)$, a naive user believing $a^\dagger \ge a_A$ perceives a higher resolution accuracy $q(a^\dagger,g) \ge q(a_A,g)$. Under Assumption~\ref{ass:reg}, Corollary~\ref{cor:mono} proves that $m^*(q)$ is weakly decreasing in accuracy, so $m^*_{\text{naive}} \le m^*_{\text{soph}}$. Furthermore, $U(m;a_A,g)$ is strictly concave in $m$ with unique maximizer $m^*_{\text{soph}}$, so any choice $m^*_{\text{naive}} \ne m^*_{\text{soph}}$ delivers strictly lower true utility: $U(m^*_{\text{naive}};a_A,g) < U(m^*_{\text{soph}};a_A,g)$.
\end{proof}

This is the source of genuine moral hazard identified so far: not raw accuracy, but
an \emph{unobserved gap between the user's belief about the assistant's policy and
the assistant's true policy}. Section~\ref{sec:model1} solves for what that true
policy $a_A$ actually is when the assistant is itself a strategic, possibly biased,
Stackelberg leader---so Proposition~\ref{prop:bias} can be instantiated with a
derived $a_A$ rather than an assumed one.

\section{The Assistant's Policy as a Stackelberg Game}
\label{sec:stackelberg}

\subsection{General formulation}

The assistant (leader) commits ex ante to a policy $\pi$ before meeting any user.
A population of users (followers) has privately known specification-cost type
$\kappa \sim F$ on $[\kappa_L,\kappa_H]$, with cost $c(m;\kappa)=\kappa m^2/2$, and
best-responds $m^*(\kappa;\pi) = \arg\max_m U(m;\pi,\kappa)$. The leader's own
objective $\Pi$ need not equal $U$: we write
\[
\Pi_\kappa(m,a) \;=\; \mu_A\, U(m,a;\kappa) \;-\; (\lambda_A - c_Q)\, a\,(k-m),
\]
where $\mu_A \in (0,1]$ discounts the leader's weight on true (downstream) user
welfare, $c_Q$ is the user's true cost of answering one clarifying question, and
$\lambda_A$ is the leader's own \emph{perceived} per-question friction cost, which
may differ from $c_Q$ (e.g., a training signal that penalizes extra turns more
than users actually mind answering a direct question). The leader is
\emph{unbiased} iff $\mu_A=1$ and $\lambda_A=c_Q$, in which case $\Pi_\kappa \equiv
U(\cdot;\kappa)$ exactly. The leader solves
$\pi^* = \arg\max_\pi \mathbb{E}_{\kappa\sim F}[\Pi_\kappa(m^*(\kappa;\pi),\pi)]$,
a bilevel program.

\begin{remark}[Existence]
For Model I below ($\pi=a\in[0,1]$, a compact scalar policy), $m^*(\kappa;a)$ is
continuous in $a$ (strict concavity of $U$ in $m$ plus the implicit function
theorem), so $\mathbb{E}_\kappa[\Pi_\kappa]$ is continuous in $a$ and a maximizer
exists on $[0,1]$ by Weierstrass's theorem. For Model II below (a finite menu over
two types), the feasible set is a compact subset of $\mathbb{R}^4$ cut out by
linear IC/IR constraints and a continuous objective, so a maximizer again exists
by Weierstrass's theorem; no ironing or ex-ante randomization is needed with only
two types.
\end{remark}

For tractability in what follows we replace the exact conjunctive success
probability $q^{k-m}$ with its first-order (small-risk) approximation around
$q\to 1$: expected loss $\approx L(1-q)(k-m)$, where $L$ denotes the economic
loss from a single unresolved-and-incorrect attribute. With $\Lambda \equiv L(1-g)$,
the user's payoff becomes
\[
U(m,a;\kappa) \;=\; V \;-\; \Lambda(1-a)(k-m) \;-\; \frac{\kappa}{2}m^2 \;-\; a\,c_Q\,(k-m).
\]

\subsection{The first-best benchmark}
\label{sec:firstbest}

Before analyzing either regime, characterize the fully-informed, fully-aligned
optimum: $(m^{FB}(\kappa), a^{FB}(\kappa)) = \arg\max_{m,a} U(m,a;\kappa)$.

\begin{proposition}[First-best is a threshold rule in $\kappa$]\label{prop:firstbest}
$U$ is linear in $a$ for fixed $m$, so $a^{FB}(\kappa) \in \{0,1\}$. Define
$\kappa^* = (\Lambda+c_Q)/(2k)$. If $\Lambda > c_Q$, then $a^{FB}(\kappa)=1$
(always ask) for $\kappa > \kappa^*$ and $a^{FB}(\kappa)=0$ (never ask, guess) for
$\kappa < \kappa^*$.
\end{proposition}

\begin{proof}
At $a=1$: $m^*(\kappa;1)=c_Q/\kappa$ and $U(m^*,1;\kappa)=V-c_Qk+c_Q^2/(2\kappa)$.
At $a=0$: $m^*(\kappa;0)=\Lambda/\kappa$ and
$U(m^*,0;\kappa)=V-\Lambda k+\Lambda^2/(2\kappa)$. Subtracting,
\[
U(a{=}1)-U(a{=}0) \;=\; [\Lambda-c_Q]\left(k - \frac{\Lambda+c_Q}{2\kappa}\right),
\]
which is positive iff $\Lambda>c_Q$ and $\kappa>(\Lambda+c_Q)/(2k)=\kappa^*$.
\end{proof}

This is a clean, closed-form, type-dependent first best: high-cost (terse) users,
who cannot cheaply self-provide detail, should always be asked; low-cost (verbose)
users, who can cheaply write enough that the small residual is safe to guess,
should never be asked. This threshold rule completely characterizes the efficient
asking policy as a function of the user's marginal cost of detail relative to downstream error loss.

\section{Model I: A Pooling (Scalar) Policy}
\label{sec:model1}

The leader is restricted to a single population-wide ask rate $a\in[0,1]$---it
cannot condition on anything user-specific. This is the natural benchmark for a
policy that is fixed at training time and cannot observe individual users.

\subsection{User best response}

The FOC $\partial U/\partial m = \Lambda(1-a) - \kappa m + ac_Q = 0$ gives, with
$\beta(\kappa)=\Lambda/\kappa$ and $\gamma(\kappa)=(\Lambda-c_Q)/\kappa$,
\[
m^*(\kappa;a) \;=\; \beta(\kappa) - a\,\gamma(\kappa), \qquad
\frac{\partial m^*}{\partial a} \;=\; -\gamma(\kappa) \;=\; \frac{c_Q-\Lambda}{\kappa}.
\]

\begin{remark}[Non-monotone effect of ask rate on specification]
Unlike Corollary~\ref{cor:mono}, the sign of $\partial m^*/\partial a$ here does
not depend on type: it is positive for \emph{every} type iff $c_Q > \Lambda$, i.e.,
iff answering a question costs users more than the expected loss it avoids. When
that holds, a higher population ask rate induces users to specify \emph{more}, not
less---they pre-empt costly clarifying questions by writing the answer down
themselves. This is a genuinely new force relative to Section~\ref{sec:baseline},
where $a$ only entered through accuracy, and is a distinct, directly testable
prediction (P1$'$ below).
\end{remark}

\subsection{The leader's problem}

Let $R(a,\kappa)=k-m^*(\kappa;a)$ denote the residual unspecified attributes. In the pooling regime, the assistant commits ex ante to an operational ask rate $a \in [0,1]$ applied population-wide, without observing or directly contracting on individual specification effort.

Evaluating the general leader objective from Section~\ref{sec:firstbest} at the follower's best response $m^*(\kappa;a) = \beta(\kappa) - a\,\gamma(\kappa)$:
\[
\Pi_\kappa(m^*(\kappa;a),a) \;=\; \mu_A\, U(m^*(\kappa;a),a;\kappa) \;-\; (\lambda_A - c_Q)\, a\,\big(k - m^*(\kappa;a)\big).
\]
Substituting $U(m^*(\kappa;a),a;\kappa) = V - \Lambda(1-a)(k-m^*) - \frac{\kappa}{2}(m^*)^2 - a\,c_Q\,(k-m^*)$, using $k - m^*(\kappa;a) = R_0(\kappa) + a\,\gamma(\kappa)$, and defining the effective resolution stake $s \equiv \Lambda - c_Q$ and the assistant's excess asking friction $b \equiv \lambda_A - c_Q$, expanding in powers of $a$ yields:
\[
\Pi_\kappa(m^*(\kappa;a),a) \;=\; \mu_A \left( V - \Lambda k + \frac{\Lambda^2}{2\kappa} \right) \;+\; (\mu_A s - b)\, R_0(\kappa)\, a \;+\; \frac{(\mu_A s - 2b)\,\gamma(\kappa)}{2}\, a^2.
\]
Averaging over $\kappa \sim F$ requires no moments beyond the population sufficient statistics already defined, $\bar R_0 = k - \mathbb{E}[\beta]$ and $\bar\gamma = \mathbb{E}[\gamma] = s\,\mathbb{E}[1/\kappa]$:
\[
\Pi(a) \;=\; \text{const} \;+\; \bar L\, a \;+\; \bar Q\, a^2,
\]
where $\text{const} = \mu_A(V - \Lambda k) + \frac{\mu_A \Lambda^2}{2}\mathbb{E}[1/\kappa]$, and
\[
\bar L \;\equiv\; (\mu_A s - b)\,\bar R_0, \qquad \bar Q \;\equiv\; \frac{(\mu_A s - 2b)\,\bar\gamma}{2}.
\]

\begin{proposition}[Stackelberg-optimal pooling rate]\label{prop:aSE}
Provided $\bar Q < 0$ (the second-order condition for strict concavity, equivalent to $2(\lambda_A - c_Q) > \mu_A(\Lambda - c_Q)$), the leader's unconstrained stationary point is
\[
a^{SE} \;=\; -\,\frac{\bar L}{2\bar Q} \;=\; -\,\frac{(\mu_A s - b)\,\bar R_0}{(\mu_A s - 2b)\,\bar\gamma}.
\]
When this stationary point lies in $[0,1]$, it is the unique global maximum of $\Pi(a)$ on $[0,1]$.
\end{proposition}

\begin{corollary}[Bias shifts the pooling rate on the interior branch]
Differentiating the closed form in Proposition~\ref{prop:aSE} directly with respect to $\lambda_A$ (holding primitives $\mu_A,\Lambda,c_Q,\bar R_0,\bar\gamma$ fixed) gives, on the interior branch $\bar Q < 0$,
\[
\frac{\partial a^{SE}}{\partial \lambda_A} \;=\; -\,\frac{\bar R_0}{\bar\gamma}\,\frac{\mu_A s}{(\mu_A s - 2b)^2} \;=\; -\,\frac{\bar R_0}{\bar\gamma}\,\frac{\mu_A (\Lambda - c_Q)}{\big[\mu_A (\Lambda - c_Q) - 2(\lambda_A - c_Q)\big]^2},
\]
which is strictly \emph{positive} whenever $\bar R_0 < 0$ (the under-specification regime). Hence, $a^{SE}$ is strictly increasing in perceived friction $\lambda_A$ throughout the interior branch. This comparative static holds on the interior branch ($\bar Q < 0$); for the global optimum in the corner regime, see Corollary~\ref{cor:corner-ray}.
\end{corollary}

\begin{corollary}[Non-identification of the two bias channels in the interior regime]
\label{cor:mu-lambda}
On the interior branch ($\bar Q < 0$), defining the normalized excess friction per unit of altruism $\tilde\rho \equiv \frac{\lambda_A - c_Q}{\mu_A} = \frac{b}{\mu_A}$, the optimal pooling rate simplifies to
\[
a^{SE} \;=\; -\,\frac{\bar R_0}{\bar\gamma} \left[ \frac{s - \tilde\rho}{s - 2\tilde\rho} \right],
\]
so $a^{SE}$ depends on the two bias parameters $(\mu_A, \lambda_A)$ strictly through the single scalar ratio $\tilde\rho = (\lambda_A - c_Q)/\mu_A$ (holding primitives $s$ and $\bar R_0/\bar\gamma$ fixed). Consequently,
\[
\frac{\partial a^{SE}/\partial\mu_A}{\partial a^{SE}/\partial\lambda_A} \;=\; -\,\frac{\lambda_A - c_Q}{\mu_A} \;=\; -\,\tilde\rho,
\]
and any two parameter pairs on the same ray $(\lambda_A - c_Q)/\mu_A = \text{const}$ produce the identical pooling ask rate. A leader that discounts downstream welfare ($\mu_A < 1$, $\lambda_A = c_Q \implies \tilde\rho = 0$) cannot be distinguished, from $a^{SE}$ alone, from any leader along the corresponding ray of excess friction.
\end{corollary}

\begin{remark}[Scope of Corollary~\ref{cor:mu-lambda}]
The derivation above is stated on the interior branch $\bar Q < 0$. By Corollary~\ref{cor:corner} below, the \emph{unbiased} pooling optimum is always a corner ($\bar Q \ge 0$), so the interior stationary point describes behavior when friction bias is sufficiently large ($2b > \mu_A s$).
\end{remark}

\begin{corollary}[Ray-invariance extends to the corner regime]\label{cor:corner-ray}
On the corner branch ($\bar Q \ge 0$, or whenever the unconstrained stationary point lies outside $[0,1]$), the optimal policy is chosen by the sign of the corner difference $\Pi(1) - \Pi(0) = \bar L + \bar Q$:
\[
\Pi(1) - \Pi(0) \;=\; \mu_A \left[ s\left(\bar R_0 + \frac{\bar\gamma}{2}\right) - \tilde\rho\,(\bar R_0 + \bar\gamma) \right].
\]
Since $\mu_A > 0$, $\mu_A$ factors out completely, and the sign depends solely on the scalar ratio $\tilde\rho = (\lambda_A - c_Q)/\mu_A$. The critical threshold ratio is
\[
\tilde\rho^{*} \;\equiv\; \frac{s\big(\bar R_0 + \bar\gamma/2\big)}{\bar R_0 + \bar\gamma},
\]
with $a^{SE} = 1$ for $\tilde\rho < \tilde\rho^{*}$ and $a^{SE} = 0$ for $\tilde\rho > \tilde\rho^{*}$. Ray-invariance in the excess friction ratio $\tilde\rho$ therefore holds unconditionally across the entire parameter space, interior and corner regimes alike. A numerical sweep of 11 rays across 10 values of $\mu_A \in [0.1, 1.0]$ (110 evaluations) confirms this analytical threshold with zero boundary-crossing violations.
\end{corollary}

\begin{corollary}[Pooling is a corner solution in the unbiased case]\label{cor:corner}
In the unbiased case ($\mu_A = 1, \lambda_A = c_Q \implies b = 0$), the quadratic coefficient is:
\[
\bar Q \;=\; \frac{s^2\,\bar\gamma}{2} \;=\; \frac{(\Lambda - c_Q)^2\,\mathbb{E}[1/\kappa]}{2} \;\ge\; 0 \quad \text{always}.
\]
The second-order condition for an interior maximum therefore fails: $\Pi(a)$ is weakly convex on $[0,1]$, its stationary point is a local minimum, and the true optimum is always a corner, $a^{SE} \in \{0,1\}$, not an interior compromise. Comparing the two corners directly:
\[
\Pi(1) - \Pi(0) \;=\; \bar L + \bar Q \;=\; s\left(\bar R_0 + \frac{\bar\gamma}{2}\right) \;=\; (\Lambda - c_Q)\left[ k - \frac{\Lambda + c_Q}{2}\,\mathbb{E}\left[\frac{1}{\kappa}\right] \right],
\]
so the unbiased pooling optimum is $a^{SE} = 1$ (always ask) when $k \ge \frac{\Lambda + c_Q}{2}\mathbb{E}[1/\kappa]$ and $a^{SE} = 0$ (never ask) when $k < \frac{\Lambda + c_Q}{2}\mathbb{E}[1/\kappa]$. This threshold matches the population-averaged first-best preference $\mathbb{E}[U(1) - U(0)]$ from Proposition~\ref{prop:firstbest} identically.
\end{corollary}

This last point matters for the comparison in Section~\ref{sec:compare}: Model I
has two independent, additive sources of loss relative to first best---the
pooling restriction, and bias.

\section{Model II: A Screening (Menu) Policy}
\label{sec:model2}

Now let the leader condition its ask rate on the user's declared specification
level (equivalently, by the revelation principle, offer a menu of bundles that
users self-select into). Consider two types $\kappa_L<\kappa_H$ with population
shares $f_L,f_H$, and let the leader choose bundles $(m_L,a_L),(m_H,a_H)$ subject
to
\[
\text{IC}_\kappa:\;\; U(m_\kappa,a_\kappa;\kappa) \ge U(m_{\kappa'},a_{\kappa'};\kappa), \qquad
\text{IR}_\kappa:\;\; U(m_\kappa,a_\kappa;\kappa) \ge \underline{U},
\]
to maximize $f_L\Pi_{\kappa_L}(m_L,a_L)+f_H\Pi_{\kappa_H}(m_H,a_H)$.

\subsection{Single-crossing}
The marginal rate of substitution between $m$ and $a$ along an indifference curve
is $\text{MRS}(\kappa) = -\partial U/\partial m \big/ \partial U/\partial a$; using
the FOC expressions above, $\partial \text{MRS}/\partial \kappa$ has a consistent
sign across the relevant range (higher-$\kappa$ types have steeper indifference
curves in $(m,a)$-space: they require more compensating ask-service per unit of
specification given up). This is the standard Spence--Mirrlees condition~\citep{spence1973} and
underwrites everything below; we omit the routine but lengthy verification.

\subsection{No distortion without bias}

\begin{proposition}[Efficiency survives private information when the leader is unbiased]
\label{prop:nodistortion}
If $\mu_A=1$ and $\lambda_A=c_Q$ (so $\Pi_\kappa \equiv U(\cdot;\kappa)$ for every
$\kappa$), then implementing each type's own first-best bundle,
$(m_\kappa,a_\kappa)=(m^{FB}(\kappa),a^{FB}(\kappa))$ from
Proposition~\ref{prop:firstbest}, is incentive compatible, and no further
mechanism design can do better for the leader.
\end{proposition}

\begin{proof}
$(m^{FB}(\kappa),a^{FB}(\kappa))$ maximizes $U(\cdot;\kappa)$ by definition, so for
any $\kappa'\ne\kappa$, $U(m^{FB}(\kappa),a^{FB}(\kappa);\kappa) \ge
U(m^{FB}(\kappa'),a^{FB}(\kappa');\kappa)$ trivially: no bundle, including another
type's, can give type $\kappa$ more utility than its own maximizer. IC holds with
no distortion needed. Since $\Pi_\kappa\equiv U(\cdot;\kappa)$, the leader's payoff
from serving each true type is exactly maximized at that type's own first best, so
this menu is also leader-optimal.
\end{proof}

This is the paper's sharpest result: unlike classical regulation (Baron--Myerson,
Laffont--Tirole), where the principal pays a literal transfer and therefore has an
independent reason to dislike conceding rent, here the leader's payoff can
coincide exactly with user welfare, and when it does, incentive compatibility is
\emph{free}---private information costs nothing. Distortion in this setting is
therefore a pure signature of misalignment, not of asymmetric information per se.

\subsection{Distortion under bias}

When $\Pi_\kappa \ne U(\cdot;\kappa)$ (bias present), the leader's preferred bundle
per type, $(m_\kappa^{B},a_\kappa^{B})=\arg\max_{m,a}\Pi_\kappa(m,a)$ (computable
in closed form by the same method as Proposition~\ref{prop:firstbest}, with
$\Lambda,c_Q$ replaced by their $\mu_A,\lambda_A$-weighted counterparts), no longer
automatically satisfies IC, because $\Pi_\kappa$ and the user's true $U(\cdot;\kappa)$
now rank bundles differently.

\begin{proposition}[Downward distortion of the high-cost type's service]
\label{prop:distortion}
Under the single-crossing condition and the empirically motivated under-asking
bias ($\lambda_A > \mu_A c_Q$, i.e.\ the leader over-weights the friction of
asking relative to true user cost), the leader's optimal incentive-compatible
menu sets $a_L^{SB}=a_L^{B}$ (no further distortion for the low-cost type, whose
biased first-best is already the binding corner $a=0$) and $a_H^{SB} < a_H^{B} \le
a_H^{FB}=1$ (the high-cost type is asked strictly less than even the leader's own
already-biased preference for them), with $a_H^{SB}$ characterized by the
standard virtual-surplus first-order condition
$\partial\Pi_{\kappa_H}/\partial a_H = \frac{f_L}{f_H}\big[\kappa_H-\kappa_L\big]
m_H \cdot \partial m_H/\partial a_H$ (rent-adjustment term on the right).
\end{proposition}

\begin{proof}[Proof sketch, corrected]
Relax the problem to IC$_L$ alone binding. Unlike classical transfer-based
screening (Baron--Myerson, Laffont--Tirole), where the principal's payoff is
directly reduced by the rent conceded to the efficient type---so the standard
technique pushes the inefficient type's IR down to reservation to minimize that
rent---here $\Pi_{\kappa_H}$ already contains $\mu_A U(\cdot;\kappa_H)$ directly,
so the leader's own payoff from serving the high-cost type \emph{rises} with
that type's true welfare: there is no independent budget motive to depress
$U_H$ to $\underline{U}$, and IR$_H$ need not bind. The bundle given to the
high-cost type is instead set by its own (biased) unconstrained optimum
$(m_H^{B},a_H^{B})=\arg\max\Pi_{\kappa_H}$, moved only as far as needed to
satisfy IC$_L$---i.e., only until $U(m_H,a_H;\kappa_L)=U(m_L,a_L;\kappa_L)$
binds---using the envelope representation of the low type's temptation to
mimic, $U(m_H,a_H;\kappa_L)-U(m_H,a_H;\kappa_H)=(\kappa_H-\kappa_L)m_H^2/2$.
This still yields $a_H^{SB}<a_H^{B}$ whenever the unconstrained biased optimum
would otherwise violate IC$_L$, but the resulting first-order condition is not
the classical virtual-surplus expression a naive import of Baron--Myerson would
suggest; see the remark below.
\end{proof}

\begin{remark}[Active constraint structure without monetary transfers]
In classical transfer-based screening (e.g., Baron--Myerson), IC$_L$ and IR$_H$ typically bind while IC$_H$ and IR$_L$ remain slack. In our setting, solving the full four-variable constrained problem shows that IC$_L$ binds alone, with IR$_H$, IC$_H$, and IR$_L$ remaining strictly slack at representative under-asking-bias parameters. This structural divergence arises because the leader's payoff $\Pi_{\kappa_H}$ directly internalizes follower welfare $\mu_A U(\cdot;\kappa_H)$, eliminating the classical principal's motive to compress agent rent down to the reservation utility. We note that the exact active constraint configuration is parameter-region-specific: under sufficiently extreme asking friction ($\lambda_A \gg c_Q$) or severe discounting ($\mu_A \ll 1$), IC$_H$ can bind alongside IC$_L$. The robust qualitative insight is that screening under misalignment strictly distorts $a_H^{SB} < a_H^B$, while the active constraint configuration reflects the non-monetary, partially aligned nature of the objective.
\end{remark}

\begin{remark}
The direction in Proposition~\ref{prop:distortion} is specific to the
under-asking bias case ($\lambda_A>\mu_A c_Q$), which we take to be the
empirically relevant one (Section~\ref{sec:naive-soph}'s motivating story: turn-level
satisfaction signals penalize extra clarifying questions more than users actually
mind answering them). Under the opposite bias, the direction of distortion can
flip; we do not claim a sign that holds for all $\Pi_\kappa$.
\end{remark}

\section{Comparing the Two Regimes}
\label{sec:compare}

\begin{corollary}[Menus weakly dominate pooling]
The leader's optimal expected payoff under Model II is weakly greater than under
Model I: any pooling policy $a$ is a feasible (degenerate) menu with
$a_L=a_H=a$, trivially satisfying IC, so the unconstrained menu optimum is at
least as good.
\end{corollary}

The more interesting comparison is \emph{where the two regimes' losses come from},
which Propositions~\ref{prop:aSE}--\ref{prop:distortion} let us decompose cleanly:

\begin{itemize}[leftmargin=*]
\item \textbf{Model I (pooling)} loses welfare from two distinct sources: (i) the pooling
restriction itself, present even when $\mu_A{=}1,\lambda_A{=}c_Q$, where it
manifests as a corner solution that fully serves one type and fully fails the
other (Corollary~\ref{cor:corner}), and (ii) bias, which shifts $a^{SE}$
(Corollary 2) and can additionally shift pooling between corner and interior
regimes. The two sources interact non-linearly across the parameter space.
\item \textbf{Model II (screening)} loses welfare from bias \emph{only}
(Proposition~\ref{prop:nodistortion}): with no bias, screening recovers full
efficiency despite private information. With bias, the loss is the
rent-extraction-driven downward distortion of Proposition~\ref{prop:distortion},
which is a genuinely different mechanism from Model I's pooling loss (it doesn't
arise from being unable to condition on anything---it arises from not being able
to trust what is conditioned on).
\end{itemize}

\subsection{Sensitivity of the Unbiased Pooling Corner}
\label{sec:unbiased_sensitivity}

While Corollary~\ref{cor:corner}'s qualitative finding---that unbiased pooling is always a corner solution rather than an interior compromise---is structural and universal, the specific corner selected depends on whether the task size $k$ exceeds the population threshold $k^* \equiv \frac{\Lambda + c_Q}{2}\mathbb{E}[1/\kappa]$. Table~\ref{tab:unbiased_sensitivity} illustrates this boundary across varying task sizes $k$ and baseline guessing accuracies $g$, holding $L=10, c_Q=2$, and $\kappa \sim \text{Uniform}[0.2, 0.6]$ (where $\mathbb{E}[1/\kappa] \approx 2.747$).

\begin{table}[ht]
\centering
\small
\begin{tabular}{cccccc}
\toprule
$g$ & $\Lambda = L(1-g)$ & Threshold $k^*$ & Task Size $k$ & $\Pi(1) - \Pi(0)$ & Optimal Corner $a^{SE}$ \\
\midrule
$0.30$ & $7.0$ & $12.36$ & $5$ & $-36.80$ & $0$ (never ask) \\
$0.30$ & $7.0$ & $12.36$ & $10$ & $-11.80$ & $0$ (never ask) \\
$0.30$ & $7.0$ & $12.36$ & $15$ & $+13.20$ & $1$ (always ask) \\
\midrule
$0.50$ & $5.0$ & $9.61$ & $5$ & $-13.84$ & $0$ (never ask) \\
$0.50$ & $5.0$ & $9.61$ & $9$ & $-1.84$ & $0$ (never ask) \\
$0.50$ & $5.0$ & $9.61$ & $10$ (Benchmark) & $+1.16$ & $1$ (always ask) \\
$0.50$ & $5.0$ & $9.61$ & $12$ & $+7.16$ & $1$ (always ask) \\
$0.50$ & $5.0$ & $9.61$ & $15$ & $+16.16$ & $1$ (always ask) \\
\midrule
$0.70$ & $3.0$ & $6.87$ & $5$ & $-1.87$ & $0$ (never ask) \\
$0.70$ & $3.0$ & $6.87$ & $8$ & $+1.13$ & $1$ (always ask) \\
$0.70$ & $3.0$ & $6.87$ & $10$ & $+3.13$ & $1$ (always ask) \\
\bottomrule
\end{tabular}
\caption{Sensitivity of the unbiased pooling corner choice $\Pi(1) - \Pi(0) = (\Lambda - c_Q)[k - k^*]$ across task complexity $k$ and guessing capability $g$, for $L=10, c_Q=2, \kappa \sim \text{Uniform}[0.2, 0.6]$. At the paper's canonical benchmark ($g=0.5, k=10$), asking yields a modest advantage of $+1.16$, which flips to never-asking at $k \le 9$.}
\label{tab:unbiased_sensitivity}
\end{table}

\begin{remark}[Compounding, testable]
Because Model II's distortion is triggered by, and scales with, the same bias
parameters ($\mu_A,\lambda_A$ vs.\ $c_Q$) that drive Model I's shift in
$a^{SE}$, reducing measured assistant bias should improve outcomes for high-cost
(terse) users under \emph{both} regimes, but the welfare \emph{gap between}
menu-conditioned and pooled policies should track population heterogeneity in
$\kappa$ (the spread $\kappa_H-\kappa_L$ and how far $F$ sits on either side of
$\kappa^*$) almost independently of bias. This gives two separable, testable
quantities from the same field data: a bias-sensitive term and a
heterogeneity-sensitive term.
\end{remark}

\section{Testable Predictions}
\label{sec:predictions}

\begin{enumerate}[leftmargin=*]
\item[\textbf{P1}] (Prop.~\ref{prop:accuracy}, Cor.~\ref{cor:mono}) Holding the
assistant's policy fixed, higher guessing accuracy should be associated with
shorter specifications but higher or equal success rates---never both shorter
specs and lower success.
\item[\textbf{P1$'$}] (Model I remark) Holding accuracy fixed, the effect of a
higher \emph{ask rate} on specification length should flip sign depending on
whether $c_Q \gtrless \Lambda=L(1-g)$: where clarifying questions are cheap
relative to the risk they resolve, users should specify less as ask rate rises
(substitution); where questions are relatively costly to answer, users should
specify \emph{more} (pre-emption). This is a within-model prediction distinct
from P1's accuracy channel.
\item[\textbf{P2}] (Prop.~\ref{prop:bias}) The correlation between turn-level
satisfaction and eventual task success should weaken, or invert, specifically for
under-asking assistants, and can now be benchmarked against the derived $a^{SE}$
of Proposition~\ref{prop:aSE} rather than an assumed ask rate.
\item[\textbf{P3}] Making ask-rate statistics visible to users should shift
specification length toward the sophisticated best response of
Section~\ref{sec:naive-soph}.
\item[\textbf{P4}] (Prop.~\ref{prop:firstbest}) Ask rates should rise with a
user's revealed difficulty specifying detail (terseness), sharply around a
threshold $\kappa^*$, rather than responding only to model confidence.
\item[\textbf{P5}] Users with lower self-reported error-catching rates should show
declining specification effort over a session (Section 3.2's dynamic channel).
\item[\textbf{P6}] (Cor.~3) Under a pooling policy, terse users should be
under-served and verbose users over-served relative to \emph{each type's own}
first best, even for a demonstrably unbiased assistant---a gap attributable
purely to instrument restriction, distinguishable in data from bias by holding
$\lambda_A$ fixed (e.g., via an A/B test on friction penalties in training) and
varying only whether the policy can condition on declared specification.
\item[\textbf{P7}] (Prop.~\ref{prop:distortion}) Even where an assistant \emph{can}
condition on declared specification, terse users should still be asked less than
the assistant's own (possibly already-biased) per-type preference would predict,
specifically because of the incentive-compatibility constraint---a distortion
that should shrink as the population's cost heterogeneity ($\kappa_H-\kappa_L$)
shrinks, holding bias fixed (Remark on compounding).
\end{enumerate}

\section{Empirical Strategy (Summary)}

The computational evaluation consists of analytical verification, numerical
experiments, and LLM-based simulation experiments; no human-subject study is
reported. Below we summarize the computational verification pipeline and outline a
proposed human-study design for future empirical research.

\textbf{Computational simulation and verification.} Instantiate LLM-simulated users with
heterogeneous cost types $\kappa$ interacting with assistants implementing Model I's
derived $a^{SE}$, Model II's numerically-solved menu, and naive baselines (fixed
confidence threshold, ignoring both stakes and type). Verify
Proposition~\ref{prop:firstbest}'s threshold numerically, verify
Proposition~\ref{prop:aSE}'s closed form against a grid search over $a$, and
verify Proposition~\ref{prop:nodistortion}/\ref{prop:distortion} by solving the
constrained menu optimization directly and checking that the low type's bundle
sits at its unconstrained optimum while the high type's does not, with the gap
shrinking as bias is set to zero.

\textbf{Proposed human-subject study design.} A between-subjects manipulation of ask
policy (pooling vs.\ menu-conditioned; biased vs.\ unbiased friction weighting) across
repeated coding tasks, measuring specification length and, where feasible,
self-reported terseness/verbosity as a proxy for $\kappa$, to test P1$'$, P6, and P7
causally rather than correlationally.

\section{Discussion and Limitations}
\label{sec:discussion}

The linear-risk approximation to the conjunctive success function is the main
tractability assumption carried through Sections~\ref{sec:stackelberg}--\ref{sec:compare}.
To verify that this simplification does not drive our core theoretical findings, we
evaluate both the unbiased pooling optimum and the biased screening menu directly under the exact conjunctive payoff $q(a,g)^{k-m}$
(independently of the linear-risk closed forms, via direct numerical grid search and
constrained optimization). We find that the corner-solution property of
Corollary~\ref{cor:corner} survives in every tested configuration across
$k\in[3,20]$ and $g\in[0.50,0.95]$. Similarly, the downward-distortion direction of
Proposition~\ref{prop:distortion} survives across the parameter grid, with exceptions concentrated at small $k$ where specification effort saturates at $m=k$ and the asking wedge collapses. Furthermore, a comprehensive parameter sweep of the full four-variable
constrained menu problem under the exact conjunctive payoff across $\lambda_A\in[c_Q,10c_Q]$ and $\mu_A\in[0.1,1.0]$ (100 grid points) confirms that $\text{IC}_H$ binds alongside $\text{IC}_L$ under extreme bias, reproducing the exact three-regime structure (no binding constraint, $\text{IC}_L$ alone, $\text{IC}_L+\text{IC}_H$ together) identified under the linear-risk approximation, with only the boundary between moderate- and extreme-bias regimes shifting slightly (from $\lambda_A\approx6.0$ to $\lambda_A\approx8.0$ at $\mu_A=1.0$). The extreme-bias reversal is therefore an intrinsic structural feature of the screening contract space, not an artifact of the linear-risk approximation.

A related limitation concerns identification of $\lambda_A$ itself. Every
bias-driven prediction (Corollary~2, Proposition~\ref{prop:distortion}, and
predictions P2, P6, P7) is stated in terms of $\lambda_A>\mu_A c_Q$, but the paper
does not currently propose a way to measure $\lambda_A$ independent of the ask-rate
behavior it is meant to explain; backing it out from observed policy alone is
circular, and Corollary~\ref{cor:mu-lambda} sharpens this into an exact
non-identification result on the interior branch: $a^{SE}$ depends on
$(\mu_A,\lambda_A)$ only through their ratio, so a discounting-driven and a
friction-driven leader with the same $\lambda_A/\mu_A$ are indistinguishable from
pooling behavior alone. An independent handle---for instance, from a training
objective or reward specification that explicitly penalizes turn
count---would be needed before the bias-versus-heterogeneity predictions in
Section~\ref{sec:predictions} can be tested non-circularly. The two-type
restriction in Model II avoids
ironing/continuum-type complications at the cost of generality; a continuum
version using the standard Mirrlees technique is the natural next step and would
let $\kappa^*$ from the first-best benchmark interact directly with the screening
distortion. We also assume truthful revelation of specified attributes (no
cheap-talk vagueness within what is said) and a single leader (no competition
among assistants, which would itself be an interesting extension---competition
plausibly disciplines bias directly, another testable channel).

A broader methodological lesson from this model concerns the hazards of importing classical contract-theoretic intuition directly into human--AI settings. In classical mechanism design (e.g., Baron--Myerson, Laffont--Tirole), the principal pays monetary transfers to extract surplus, so agent rent directly diminishes principal payoff. In human--AI interaction, the assistant's objective directly incorporates (a discounted multiple of) user welfare. As a result, incentive compatibility can be free in the absence of bias (Proposition~\ref{prop:nodistortion}), participation constraints do not automatically bind at the bottom of the type distribution, and unbiased pooling naturally induces corner solutions rather than interior compromises (Corollary~\ref{cor:corner}). Contract-theoretic treatments of AI interaction must therefore explicitly account for the non-monetary, partially aligned nature of the interaction rather than relying on analogies to cash-transfer mechanisms.

\bibliographystyle{plainnat}
\begin{thebibliography}{13}

\bibitem[von Stackelberg(1934)]{stackelberg1934}
Heinrich von Stackelberg. \textit{Marktform und Gleichgewicht}. Julius Springer, Berlin/Wien, 1934.

\bibitem[Myerson(1981)]{myerson1981}
Roger B.\ Myerson. Optimal auction design. \textit{Mathematics of Operations
Research}, 6(1):58--73, 1981.

\bibitem[Baron and Myerson(1982)]{baron-myerson1982}
David P.\ Baron and Roger B.\ Myerson. Regulating a monopolist with unknown
costs. \textit{Econometrica}, 50(4):911--930, 1982.

\bibitem[Mussa and Rosen(1978)]{mussa-rosen1978}
Michael Mussa and Sherwin Rosen. Monopoly and product quality. \textit{Journal
of Economic Theory}, 18(2):301--317, 1978.

\bibitem[Laffont and Tirole(1986)]{laffont-tirole1986}
Jean-Jacques Laffont and Jean Tirole. Using cost observation to regulate firms.
\textit{Journal of Political Economy}, 94(3):614--641, 1986.

\bibitem[Crawford and Sobel(1982)]{crawford-sobel1982}
Vincent P.\ Crawford and Joel Sobel. Strategic information transmission.
\textit{Econometrica}, 50(6):1431--1451, 1982.

\bibitem[Townsend(1979)]{townsend1979}
Robert M.\ Townsend. Optimal contracts and competitive markets with costly state
verification. \textit{Journal of Economic Theory}, 21(2):265--293, 1979.

\bibitem[Grossman(1981)]{grossman1981}
Sanford J.\ Grossman. The informational role of warranties and private
disclosure about product quality. \textit{Journal of Law and Economics},
24(3):461--483, 1981.

\bibitem[Sims(2003)]{sims2003}
Christopher A.\ Sims. Implications of rational inattention. \textit{Journal of
Monetary Economics}, 50(3):665--690, 2003.

\bibitem[Spence(1973)]{spence1973}
Michael Spence. Job market signaling. \textit{Quarterly Journal of Economics},
87(3):355--374, 1973.

\bibitem[Topkis(1978)]{topkis1978}
Donald M.\ Topkis. Minimizing a submodular function on a lattice.
\textit{Operations Research}, 26(2):305--321, 1978.

\bibitem[Topkis(1998)]{topkis1998}
Donald M.\ Topkis. \textit{Supermodularity and Complementarity}.
Princeton University Press, Princeton, NJ, 1998.

\bibitem[Milgrom and Shannon(1994)]{milgrom-shannon1994}
Paul Milgrom and Christina Shannon. Monotone comparative statics.
\textit{Econometrica}, 62(1):157--180, 1994.

\end{thebibliography}

\end{document}
